\begin{abstract}
Every firm that trades keeps score more than once. One set of systems books the
trades; another values them; a third reports them; a fourth settles them. These
systems hold overlapping copies of the same facts, and the copies drift apart.
When they disagree --- and they disagree every day --- people reconcile: they
find the breaks, decide which copy is right, and correct the others by hand. This
document describes a system of record for post-trade activity built so that a
large class of those disagreements cannot occur, because the numbers they would
disagree about are never stored twice. Positions, profit and loss, balance
sheets, and reports are computed from one immutable record rather than maintained
alongside it.

The design rests on a single claim. \emph{The ledger is the orchestrator of the
output of many smart contracts.} The smart contracts run outside the ledger's
trust boundary, each triggered by an event --- an incoming trade, a lifecycle
event (an expiry, a coupon, a barrier fixing), or a corporate action (a dividend,
a split, a spin-off). When its trigger fires, a contract reads the declared data
in the ledger and the stamped observations it is allowed to see, works out what
must change, and produces transactions --- moves, terms versions, status writes.
The ledger runs none of the contracts' product-specific logic. Transactions
arrive at its single door; the door admits each one, records it in the immutable
log, and every view --- positions, profit and loss, reports --- is computed from
the log.

The chapters that follow substantiate this claim: first the door and the log,
then the declared data a contract may read, then the contracts and their
triggers, and finally a sequence of real products --- cash equity, a dividend
forecast, listed futures, a structured note, a special dividend, a spin-off, an
elective merger, and a securities loan --- each processed by the same machine.
\end{abstract}

\tableofcontents
\clearpage

\section{The Orchestrator Thesis}\label{sec:thesis}

\subsection{The problem is reconciliation}

Consider what has to be true for two internal reports to agree. A trading desk's
risk system says the book holds 2{,}000 shares of a company; the firm's general
ledger says the same position is worth \EUR{100{,}000}; the regulatory report
says a dividend was received on it last week. These three statements are about one
underlying fact --- a holding and its history --- but each system stores its own
copy of that fact and updates it on its own schedule. If the risk system applied
a stock split and the general ledger did not, the share count and the valuation
now describe different states of the world. Nothing inside either system is wrong;
they are simply out of step. Finding that break, and fixing it, is reconciliation,
and it is where much of the cost and most of the operational risk of a
post-trade system lives.

The design in this document removes the source of the break rather than
detecting it faster. There is one record: an ordered, append-only stream of
events. A position is not a stored number that systems update; it is the sum of
the moves that touched it, computed on demand from the stream. Profit and loss is
not a second stored number; it is the difference of two valuations, each computed
from the same positions. A balance sheet is a projection; a report is a
projection. Because balances, profit and loss, and reports are all computed from
the one stream, two of them cannot disagree about a fact the stream records only
once. Within the system's scope, three reconciliation failures that dominate
real operations --- trading system against general ledger, profit and loss
against balance sheet, and one entity's books against another's --- are not
merely made rare. Each of them is a disagreement between two projections of one
record, and a disagreement between two projections of one record is not a state
the system can be in.

\subsection{The division of labour}

The stream is not enough on its own, because someone has to decide what events to
put in it. A trade must become a change in holdings and cash; a coupon date must
become a payment; a two-for-one split must double a share count and halve a price
reference. Each of these decisions is specific to a product and, often, to a
single contract's terms. There is no general rule that turns ``a barrier was
observed at 78.40'' into the right set of moves for every structured note ever
written; the note's own terms say what happens.

The design draws a sharp line here. The logic that works out what a particular
event means for a particular product lives in a \emph{smart contract}: a
deterministic function that reads the relevant facts and produces the resulting
transactions. A financial instrument, on this view, is a contract that emits
moves, and adding a new product is writing a new contract, not changing the
ledger. The ledger itself performs none of this product-specific work. It admits
transactions; it does not manufacture them. Its job is to take each proposed
transaction at its single door, decide whether the transaction is well-formed,
record it if so, and refuse it otherwise --- and then to answer every question
about the resulting state by reading the log.

This division is what the word \emph{orchestrator} in the thesis names. The
contracts compute; the ledger admits and records; the views are read back. The
ledger orchestrates the output of the contracts without running the contracts'
internals, in the same sense that a database commits the rows an application hands
it without containing the application's business rules.

\subsection{Three words that need pinning down}

The thesis uses three phrases loosely enough to mislead a careful reader, so we
fix their exact content before building on them.

\paragraph{``Outside the ledger'' means outside its trust boundary.} A contract
is external to the ledger, untrusted by it, and verifiable against it. None of
this requires the contract to run in a separate process or on a separate machine.
The reference implementation of a processor is an ordinary function called in the
same process as the door. ``Outside'' is a statement about authority, not
location: the ledger never treats a contract's output as authoritative merely
because the contract produced it. Whatever a contract submits is a proposal until
the door admits it, and it is admitted under exactly the same checks as every
other transaction.

\paragraph{``Runs none of the logic'' is scoped to product-specific logic.} The
ledger interprets nothing and prices nothing. It runs no pricing model, no
product-specific settlement rule, no event-specific special case. It does,
however, run one generic thing, and honesty requires naming it. When a datum
recorded before a corporate action is read after it, the ledger converts that
datum by an operator the corporate action \emph{declared} --- halving a price
across a two-for-one split, for instance. This conversion is not product logic
and not a pricing decision; it is the mechanical application of a declared
operator drawn from a small fixed menu, and the ledger applies it uniformly to
everything that crosses the boundary. The pricing function that decided a note was
worth what it was worth stays a black box the ledger never opens. Section~\ref{sec:data}
makes this precise; here it is enough to say that ``the ledger runs none of the
logic'' means none of the \emph{contracts'} logic. The one computation the ledger
does perform is this read-seam application, and it is not an interpreter of anything
a contract wrote.

\paragraph{Conservation is not something the door checks.} It would be natural to
say the door admits a transaction only after checking that it does not create or
destroy value. That is not what happens, and the difference matters. The unit of
change is a \emph{move}: a positive quantity of one unit leaves one wallet and the
same quantity enters another. A move is a transfer and nothing else; it has no way
to express a quantity appearing from nowhere. Because a transaction is built from
moves, an unconserved transaction cannot be constructed --- there is no value for
the door to reject. The door checks the things that \emph{can} fail. Conservation
is not among them, because it cannot fail. Section~\ref{sec:door} develops this;
it is the first and strongest instance of a principle the whole design leans on:
where an illegal state can be made unrepresentable, that is preferred to checking
for it, because a check can be forgotten and a type cannot.

\subsection{Trusted for nothing, verifiable in everything}

If the contracts are untrusted, why is it safe to record what they produce?
Because everything a contract reads is in the log. A contract reads declared
terms, the adjustment schedule, dimensions, the datum-kind registry, and stamped
observations --- all of them log-resident. Given the log up to the point where a
contract ran, anyone can run the same computation again and check that the
contract's output is what the declared inputs imply. The ledger therefore trusts
a contract to \emph{produce} a transaction, never to be \emph{right} about it;
being right is not trusted, it is checked. This is what ``trusted for nothing,
verifiable in everything'' asserts, and for everything a contract produces the
assertion is unconditional.

For what the system takes in from the outside world, the assertion is not
unconditional, and it would be dishonest to claim otherwise. Two inputs rest on
named, bounded trust assumptions. First, when a market observation is stamped on
the way in, the stamp is trusted to describe the source's dissemination
truthfully. Second, before a kind of observation may be consumed at all, it is
registered, and the registration declares how that kind transforms across a
boundary. That declaration is trusted --- but only within a bound. A mandatory
per-entry witness checks the declaration against the value-invariance identity and
refuses any kind whose declared dimension would not survive a scale; the trust
reaches exactly as far as that check and no further. These two assumptions are
stated once, carried explicitly, and made to fail closed. They are not eliminated. A further mechanism illustrates the limit, and it is a
rule of the ledger, not a judgement about the market. Every stamped datum declares
the basis it was observed in. If a datum asserts a basis for which no boundary has
been recorded --- it claims to live one frame forward of anything the log knows ---
the ingest door refuses it and quarantines the series by policy. The rule fires on
a structural mismatch between the datum's declared basis and the log, never on an
economic guess that some price move ``must'' have been a corporate action. What
lies outside the system is only the human follow-up: obtaining the missing attested
notice, so the boundary can be recorded and the series released. That is detection
and repair after the fact, not prevention, and the system is honest about the
difference.

\subsection{Two layers of correctness}

The distinction between what is guaranteed and what is checked runs through the
whole design, and it is worth stating once, plainly, as two layers.

\begin{principle}[Admission and recomputation]\label{prin:two-layer}
Admission guarantees that no transaction can break the ledger; recomputability
from declared data guarantees that a wrong transaction can always be detected.
Structural correctness is enforced at the door. Economic correctness is not
enforced at the door: it is made checkable by recomputing the transaction from the
log, and a failed check marks a producer defect, repaired downstream.
\end{principle}

The first layer is unconditional and owes nothing to the contract. A transaction
that reaches the log conserves every unit, applies in full or not at all, writes
each fact through its one legal writer, and leaves no dangling reference --- and
it has these properties whether the contract that produced it was careful,
careless, or hostile, because the properties never depended on the contract. The
second layer is where a contract can still be wrong: it can propose a transaction
that is perfectly well-formed and yet books the wrong economics --- pays
\EUR{98} where the declared terms owe \EUR{100}. The door has no basis to refuse
such a transaction, and it does not try. Instead the discrepancy is caught by
recomputing the transaction from the log and comparing; the difference is exactly
the defect, and it is a defect in the producer, never in the ledger, repaired by
a compensating transaction. The ledger does not manufacture that repair. The
producer is corrected and resubmits through the same door, and the resubmission
carries a \texttt{corrects} field naming the transaction it supersedes; the
original stays in the log, now marked corrected, and both the error and its repair
survive. Section~\ref{sec:contracts} works this example through in full.

The two layers reduce to one sentence: the ledger \emph{relies on} the first and
\emph{expects} the second. It relies on
structural correctness because structural correctness is guaranteed by
construction and no producer can break it. It expects economic correctness of its
contracts, and rather than trusting that expectation it makes it checkable,
turning ``the contract must be right'' from an article of faith into a
recomputation anyone can run.

The remainder of the document is this claim, developed and then substantiated.
Section~\ref{sec:door} builds the door and the log and shows why an untrusted
proposal is safe to record. Section~\ref{sec:data} defines the declared data a
contract may read. Section~\ref{sec:contracts} specifies the contracts and proves
out the two layers. The product tour that follows runs eight economically
distinct events through the one machine, and each stop closes one design question
by showing the mechanism absorb a real product with real numbers.
